\scanpage{27}
contenant les $x_i$, et pour un idéal $K$ stable \inserted{(et pour un tel seulement)}, les $\pi^n$ passent au quotient et définissent sur $(A/K,J/K)$ une structure d'anneau à idéal avec puiss. divisées.]

Anneau avec idéal à puissances divisées \uline{$n$-tronqué}: qui satisfait
\[
\deleted{\unsure{(x)}}\quad x^{(i)}=0\quad\text{pour }i>\unsure{n|m},\ x\in J
\]
ce qui implique que
\[
x\in J\Longrightarrow x^{n+1}=0
\]
Si $J$ est engendré par un nb fini d'éléments $x_1---x_N$, tous ceux-ci sont nilpotents, $J$ est un \uline{nilidéal} [de façon précise, $J^{Nn+1}=0$]. \deleted{Application}

Anneau ``avec'' \uline{nilidéal à puissances divisées \deleted{``nilpotentes}} \deleted{$\forall x,\exists n$} $\exists n$, tel \deleted{\unsure{que}} qu'il soit \deleted{tronqué} à puissances divisées $n$-tronquées. Anneau ``avec idéal à puissances divisées nilpotentes'', $\forall x$, $\exists n$ tel que $x^{(i)}=0$ si $i>n$. On a, pour un tel anneau, deux isom de groupes inverses l'un de l'autre

\diagram{\begin{tikzcd}[ampersand replacement=\&,column sep=huge,row sep=large]
1+J \arrow[r,shift left=1.0ex,"\log(1+x)=\sum_{\unsure{n>0|n\geq0}}(-1)^n(n-1)!x^{(n)}"] \& J \arrow[l,shift left=1.0ex,"\exp x=\sum_{n\geq0}x^{(n)}"] \\
A^* \arrow[u,phantom,"\cap" description] \&
\end{tikzcd}}
\mspage{26}
